%%% BEGIN fig:yzliu-complex-lorentz
\begin{figure}[htbp]
\centering
\begin{tikzpicture}[fd,scale=1.08,
  fd2 ax/.style={lcaxis,gray!60!black},
  fd2 lc/.style={draw=gray!55,line width=0.5pt,dash pattern=on 3pt off 2pt},
  fd2 wl/.style={wilson=#1},
  fd2 wlc/.style={draw=violet!70!black,solid,double,double distance=1.6pt,line width=0.55pt,
      postaction={decorate,solid,decoration={markings,mark=at position #1 with
      {\arrow[xshift=3.2pt,line width=0.6pt,violet!70!black]{Stealth[length=6.4pt,width=6pt]}}}}},
  fd2 vtx/.style={circle,fill=black,draw=none,inner sep=0pt,minimum size=4pt},
  fd2 vtxc/.style={circle,fill=violet!70!black,draw=none,inner sep=0pt,minimum size=4pt},
  fd2 shift/.style={draw=#1,line width=0.6pt,-{Stealth[length=4pt,width=3.4pt]}},
  fd2 lab/.style={font=\small},
  fd2 note/.style={font=\footnotesize}]
  % ===================== left: S, DY shape, exponential regulator =====================
  \begin{scope}
    \draw[fd2 ax] (-2.15,1.0) -- (-1.55,1.0) node[right=-1pt,fd2 note] {$z$};
    \draw[fd2 ax] (-2.15,1.0) -- (-2.15,1.6) node[above=-1pt,fd2 note] {$t$};
    \draw[fd2 lc] (-1.85,-1.85) -- (1.25,1.25);
    \draw[fd2 lc] (1.85,-1.85) -- (-1.25,1.25);
    % amplitude T[W^dag_nbar(0) W_n(0)]: past-pointing lines meeting at 0
    \draw[fd2 wl=0.55] (-1.5,-1.5) -- (0,0);
    \draw[fd2 wl=0.72] (0,0) -- (1.5,-1.5);
    \node[fd2 vtx] at (0,0) {};
    % conjugate amplitude Tbar[W^dag_n W_nbar] at b - i nu e_t (drawn displaced)
    \draw[fd2 wlc=0.55] (1.5,-2.65) -- (0,-1.15);
    \draw[fd2 wlc=0.55] (0,-1.15) -- (-1.5,-2.65);
    \node[fd2 vtxc] at (0,-1.15) {};
    \draw[fd2 shift=violet!60!black,dotted,line width=0.8pt] (0,-0.1) -- (0,-1.04);
    \node[font=\scriptsize,text=violet!55!black,anchor=west,inner sep=1pt] at (0.04,-0.9) {$-\ii\nu e_t$};
    % labels
    \node[fd2 note,anchor=west] at (0.1,0.02) {$0$};
    \node[fd2 note,anchor=north east] at (-1.42,-1.45) {$n$};
    \node[fd2 note,anchor=north west] at (1.42,-1.45) {$\bar n$};
    \node[fd2 note,anchor=south] at (0,0.14) {$T[W_{\bar n}^\dagger W_n](0)$};
    \node[fd2 note,anchor=north,align=center,text=violet!55!black] at (0,-2.75) {$\bar T[W_n^\dagger W_{\bar n}](\vec b_\perp-\ii\nu e_t)$\\[1pt]（虚时间位移）};
    \node[fd2 lab,align=center] at (0,-4.2) {$S$：DY 形软因子（指数正规子）};
  \end{scope}
  % ===================== right: S_t, single time ordering, space-like form factor =====================
  \begin{scope}[xshift=9.6cm]
    \draw[fd2 ax] (0.55,1.0) -- (1.15,1.0) node[right=-1pt,fd2 note] {$z$};
    \draw[fd2 ax] (0.55,1.0) -- (0.55,1.6) node[above=-1pt,fd2 note] {$t$};
    \draw[fd2 lc] (-1.85,-1.85) -- (1.0,1.0);
    \draw[fd2 lc] (1.0,-1.0) -- (-1.85,1.85);
    % pair at 0: incoming along n (past), outgoing along nbar (future)
    \draw[fd2 wl=0.55] (-1.5,-1.5) -- (0,0);
    \draw[fd2 wl=0.6] (0,0) -- (-1.5,1.5);
    \node[fd2 vtx] at (0,0) {};
    % pair at b - nu e_z (real space-like shift)
    \draw[fd2 wlc=0.5] (-2.7,1.5) -- (-1.2,0);
    \draw[fd2 wlc=0.55] (-1.2,0) -- (-2.7,-1.5);
    \node[fd2 vtxc] at (-1.2,0) {};
    \draw[fd2 shift=violet!60!black] (-0.1,0) -- (-1.1,0);
    \node[font=\scriptsize,text=violet!55!black,anchor=south,inner sep=1pt] at (-0.68,0.05) {$-\nu e_z$};
    % labels
    \node[fd2 note,anchor=west] at (0.1,0) {$0$};
    \node[fd2 note,anchor=north west] at (-1.42,-1.45) {$n$};
    \node[fd2 note,anchor=south west] at (-1.42,1.45) {$\bar n$};
    \node[fd2 note,anchor=north,align=center,text=violet!55!black] at (-2.0,-1.62) {$\vec b_\perp-\nu e_z$\\[1pt]（实的类空位移）};
    \node[fd2 note,anchor=south,align=center] at (-1.2,1.95) {$T[W_n^\dagger W_{\bar n+}^\dagger\,W_{\bar n+}W_n]$（单一时序）};
    \node[fd2 lab,align=center] at (-1.1,-4.2) {$S_t$：单一时序的类空形状因子};
  \end{scope}
  % ===================== the complex Lorentz transformation =====================
  \draw[line width=1.5pt,{Stealth[length=8pt,width=8pt]}-{Stealth[length=8pt,width=8pt]}] (2.25,0.05) -- (6.3,0.05);
  \node[fd2 lab,anchor=south] at (4.27,0.22) {$\Lambda(t,z)=(\ii z,\ii t)$};
  \node[fd2 note,anchor=north,align=left] at (4.27,-0.2) {$\Lambda(-\ii\lambda n)=\lambda n$\\[1pt]$\Lambda(-\ii\lambda\bar n)=-\lambda\bar n$\\[1pt]$\Lambda(-\nu e_z)=-\ii\nu e_t$};
  % ===================== invariant lengths =====================
  \node[fd2 lab] at (4.3,-5.0) {$x^2_{B,ij}=-\nu^2-b_\perp^2-\sqrt2\,\ii\nu\big(\lambda^L_i-\lambda^R_j\big)$：实部 $\le-b_\perp^2-\nu^2<0$，阈下、无割线};
\end{tikzpicture}
\caption{时序无关性的解析延拓图像（$(z,t)$ 平面示意；双线为类光 Wilson 线，箭头为路径方向；$n$、$\bar n$ 分别沿 $+z$、$-z$，即本书 $n_A$、$n_B$ 的方向；横向间隔 $\vec b_\perp$ 垂直于纸面）。
左：指数正规子下的 DY 形软因子 $S$——振幅中的 $T[W_{\bar n}^\dagger W_n](0)$ 与共轭振幅中位于 $\vec b_\perp-\ii\nu e_t$ 的 $\bar T[W_n^\dagger W_{\bar n}]$（紫色，虚时间位移画成竖直错开）；
右：单一时序的类空形状因子 $S_t$——位于 $0$ 与 $\vec b_\perp-\nu e_z$ 的一对类光色源，沿 $n$ 入射、沿 $\bar n$ 出射。
二者由复 Lorentz 变换 $\Lambda(t,z)=(\ii z,\ii t)$（快度为 $\ii\pi/2$ 的推动）相联系，所有不变间隔的实部都为负（阈下表示，例如图下的 $x^2_{B,ij}$）。
在维数正规化中胶子 Wightman 函数满足谱条件与局域对易性的假设下（\cref{asm:yzliu-wightman-pt}；复 Lorentz 协变性来自 Bargmann--Hall--Wightman 定理），微扰论各阶 $S=S_t$（\cref{thm:yzliu-time-ordering}，全阶论证 [A]；VF-YZL-12）。}
\label{fig:yzliu-complex-lorentz}
\end{figure}
%%% END fig:yzliu-complex-lorentz

%%% BEGIN fig:yzliu-leading-region
\begin{figure}[htbp]
\centering
\begin{tikzpicture}[fd,scale=0.95,
  fd2 soft/.style={gluon,draw=green!45!black},
  fd2 hadron/.style={bound state,line cap=butt},
  fd2 dipole/.style={draw=green!40!black,line width=2.2pt,dash pattern=on 5pt off 3pt,line cap=butt},
  fd2 lab/.style={font=\small},
  fd2 note/.style={font=\footnotesize}]
  % ======================= (a) double-line notation =======================
  \begin{scope}
    % staple (space-like): from psi(0) along +z to L, transverse link, back along -z to psibar
    \draw[wilson=0.5] (1.94,1.75) -- (4.3,1.75);
    \draw[wilson=0.55] (4.3,1.75) -- (4.3,2.8);
    \draw[wilson=0.42] (4.3,2.8) -- (0.69,2.8);
    % soft subgraph S: gluons to the staple and to C
    \draw[fd2 soft] (2.98,1.07) -- (2.55,1.75);
    \draw[fd2 soft] (3.55,1.05) -- (4.3,2.25);
    \draw[fd2 soft] (2.85,0.55) -- (2.4,0.3);
    \draw[fd2 soft] (3.1,0.3) -- (2.6,0.0);
    \node[soft blob,circle,minimum size=0.9cm,inner sep=0pt] at (3.35,0.68) {$S$};
    % collinear quark lines (arrows: fermion-number flow) and longitudinal gluons
    \draw[quark=0.55] (1.68,0.48) -- (1.7,1.47);
    \draw[quark=0.5] (0.45,2.52) -- (0.55,0.5);
    \draw[longit] (2.12,0.42) -- (1.88,1.55);   \longitmark{(2.12,0.42)}{(1.88,1.55)}
    \draw[longit] (-0.02,0.38) -- (0.26,2.58);  \longitmark{(-0.02,0.38)}{(0.26,2.58)}
    \node[fd2 note,text=blue!55!black] at (2.33,1.05) {$\cdots$};
    \node[fd2 note,text=blue!55!black] at (-0.27,1.5) {$\cdots$};
    % collinear subgraph C (hadron P along +z; forward matrix element)
    \draw[fd2 hadron] (-0.95,-0.75) -- (0.25,-0.2);
    \draw[fd2 hadron] (2.55,-0.2) -- (3.75,-0.75);
    \node[jet blob,ellipse,minimum width=2.7cm,minimum height=1.05cm,inner sep=0pt] at (1.4,0.0) {$C$};
    % hard vertices at the two quark--link junctions
    \node[hard blob,circle,minimum size=0.54cm,inner sep=0pt] at (1.7,1.75) {\footnotesize$H$};
    \node[hard blob,circle,minimum size=0.54cm,inner sep=0pt] at (0.45,2.8) {\footnotesize$H$};
    % labels
    \node[fd2 note,anchor=north west] at (-0.95,-0.75) {$P$};
    \node[fd2 note,anchor=north east] at (3.75,-0.75) {$P$};
    \node[fd2 note,anchor=east] at (1.4,1.78) {$0$};
    \node[fd2 note,anchor=south] at (0.45,3.08) {$z\hat e_z+\vec b_\perp$};
    \node[fd2 note,anchor=west] at (4.38,2.27) {$L\to\infty$};
    \node[fd2 lab] at (1.9,-1.5) {(a)\ 双线记号};
  \end{scope}
  % ======================= (b) single-line notation =======================
  \begin{scope}[xshift=7.15cm]
    % dipole line H -- S and soft gluons S -- C
    \draw[fd2 dipole] (1.33,2.3) -- (3.12,2.3);
    \draw[fd2 soft] (3.35,1.86) -- (2.42,0.3);
    \draw[fd2 soft] (3.8,1.9) -- (2.66,0.02);
    \node[soft blob,circle,minimum size=0.9cm,inner sep=0pt] at (3.55,2.3) {$S$};
    % collinear quark lines and longitudinal gluons into H
    \draw[quark=0.55] (0.75,0.48) -- (0.75,1.97);
    \draw[quark=0.55] (1.15,1.97) -- (1.2,0.5);
    \draw[longit] (1.75,0.42) -- (1.34,2.07);   \longitmark{(1.75,0.42)}{(1.34,2.07)}
    \draw[longit] (0.12,0.4) -- (0.56,2.07);    \longitmark{(0.12,0.4)}{(0.56,2.07)}
    % C and H
    \draw[fd2 hadron] (-0.95,-0.75) -- (0.25,-0.2);
    \draw[fd2 hadron] (2.55,-0.2) -- (3.75,-0.75);
    \node[jet blob,ellipse,minimum width=2.7cm,minimum height=1.05cm,inner sep=0pt] at (1.4,0.0) {$C$};
    \node[hard blob,circle,minimum size=0.75cm,inner sep=0pt] at (0.95,2.3) {$H$};
    \node[fd2 note,anchor=north west] at (-0.95,-0.75) {$P$};
    \node[fd2 note,anchor=north east] at (3.75,-0.75) {$P$};
    \node[fd2 note,anchor=south,text=green!30!black] at (2.25,2.82) {偶极线：量纲 $\frac32$、自旋 $0$};
    \node[fd2 lab] at (1.4,-1.5) {(b)\ 单线记号};
    % power counting
    \node[draw=gray!60,rounded corners=2pt,fill=gray!5,inner sep=4pt,align=left,font=\footnotesize,anchor=west] at (4.35,0.95)
      {$F(H)=-\frac72+2+\frac32=0$\\[3pt]
       $2$：两条共线夸克线\\[1pt]
       $\frac32$：一条偶极线\\[3pt]
       $F(SC)\ge0$};
  \end{scope}
\end{tikzpicture}
\caption{准 TMD 的领头区域（示意；幂次计数参数为 $1/(b_\perp P^z)$）。
(a)~双线记号：共线子图 $C$（强子动量 $P$ 沿 $+\hat z$，前向矩阵元）、位于两个夸克—规范线连接点 $0$ 与 $z\hat e_z+\vec b_\perp$ 的硬顶点 $H$、类空订书钉（双线，箭头为路径方向：自 $\psi(0)$ 沿 $+\hat z$ 到 $L\to\infty$，经横向连接回到 $\bar\psi$；与\cref{fig:tmd-staple}中的类光订书钉不同）与软子图 $S$；细线箭头为费米子数流向；$C$ 与 $H$ 之间可有任意多条纵向共线胶子（带点的蓝色胶子），$S$ 与 $C$ 之间可有任意多条软胶子（绿色）。
(b)~单线记号：一对订书钉线合成一条量纲 $\frac32$、自旋 $0$ 的偶极线（粗虚线）。
幂次计数 $F(H)=0$、$F(SC)\ge0$ 表明没有超领头区域，且硬动量不能在 $0$ 与 $\vec b_\perp$ 之间传递，因而匹配中没有 $x$ 卷积（VF-YZL-8；\cref{prop:yzliu-qtmd-powercount}，全阶论证 [A]）。}
\label{fig:yzliu-leading-region}
\end{figure}
%%% END fig:yzliu-leading-region

%%% BEGIN fig:yzliu-soft-hqet
\begin{figure}[htbp]
\centering
\begin{tikzpicture}[fd,
  fd2 ax/.style={lcaxis,gray!60!black},
  fd2 lc/.style={draw=gray!55,line width=0.5pt,dash pattern=on 1.2pt off 1.6pt},
  fd2 warrow/.style={postaction={decorate,decoration={markings,
      mark=at position #1 with {\arrow[xshift=3.2pt]{Stealth[length=6.4pt,width=6pt]}}}}},
  fd2 heavy/.style={draw,solid,line width=1.4pt,fd arrow=#1},
  fd2 heavy/.default=0.5,
  fd2 sheet/.style={fill=gray!13,draw=none},
  fd2 lab/.style={font=\small},
  fd2 note/.style={font=\footnotesize}]
  % geometry: beta = 0.55, t = t' = 1.85; perspective b_perp vector B = (1.7,0.55)
  \foreach \panel/\xs in {a/0, b/7.6} {
  \begin{scope}[xshift=\xs cm]
    \coordinate (O) at (0,0);          \coordinate (OB) at (1.7,0.55);
    \coordinate (P) at (-1.0175,-1.85); \coordinate (PB) at (0.6825,-1.3);
    \coordinate (F) at (-1.0175,1.85);  \coordinate (FB) at (0.6825,2.4);
    % axes and light cone (b_perp = 0)
    \draw[fd2 lc] (-1.75,-1.75) -- (1.75,1.75);  \draw[fd2 lc] (-1.75,1.75) -- (1.75,-1.75);
    \draw[fd2 ax] (-1.75,0) -- (2.75,0) node[below left=1pt and -3pt,font=\footnotesize] {$z$};
    \draw[fd2 ax] (0,-2.25) -- (0,2.95) node[left,font=\footnotesize] {$t$};
    \draw[gray!60!black,line width=0.5pt,dash pattern=on 3pt off 2pt,-{Stealth[length=4pt,width=3.4pt]}] (O) -- (2.45,0.79) node[right,font=\footnotesize] {$\vec b_\perp$};
    \if a\panel
      % two staple-shaped Wilson lines: past (v) and future (v'), closed at -t and t'
      \draw[wilson plain] (OB) -- (PB) -- (P) -- (O) -- (F) -- (FB) -- cycle;
      \path[fd2 warrow=0.55] (P) -- (O);    \path[fd2 warrow=0.6] (O) -- (F);
      \path[fd2 warrow=0.55] (FB) -- (OB);  \path[fd2 warrow=0.6] (OB) -- (PB);
      \path[fd2 warrow=0.55] (F) -- (FB);   \path[fd2 warrow=0.55] (PB) -- (P);
      \node[opvertex] at (O) {};  \node[opvertex] at (OB) {};
      \node[fd2 lab,anchor=east] at (-0.6,-1.0) {$v$};
      \node[fd2 lab,anchor=east] at (-0.6,1.0) {$v'$};
      \node[fd2 note,anchor=north east] at ($(P)+(0.02,-0.04)$) {$-t$};
      \node[fd2 note,anchor=south east] at ($(F)+(0.02,0.04)$) {$t'$};
      \node[fd2 note,anchor=north west] at (0.05,-0.02) {$0$};
      \node[fd2 lab,align=center] at (0.45,-3.1) {$S=\displaystyle\lim_{t,t'\to\infty}\frac{W(t,t',b_\perp,Y,Y')}{\sqrt{Z(2t,b_\perp,Y)\,Z(2t',b_\perp,Y')}}$};
      \node[fd2 note,align=left,anchor=west] at (1.95,1.75) {时序乘积 $\mathcal T$};
      \node[fd2 lab] at (0.45,-4.2) {(a)\ 订书钉形 Wilson 回路};
    \else
      % HQET picture: static colour sources Q (at 0) and Qbar (at b_perp)
      \fill[fd2 sheet] (P) -- (O) -- (F) -- (FB) -- (OB) -- (PB) -- cycle;
      \draw[fd2 heavy=0.55] (P) -- (O);    \draw[fd2 heavy=0.55] (O) -- (F);
      \draw[fd2 heavy=0.5] (FB) -- (OB);   \draw[fd2 heavy=0.55] (OB) -- (PB);
      \draw[gray!60!black,line width=0.6pt,dash pattern=on 2.5pt off 1.5pt] (O) -- (OB);
      \node[opvertex] at (O) {};  \node[opvertex] at (OB) {};
      \node[fd2 lab,anchor=east] at (-0.6,-1.0) {$v$};
      \node[fd2 lab,anchor=east] at (-0.6,1.0) {$v'$};
      \node[fd2 note,anchor=east] at (-1.05,-1.9) {$Q$};
      \node[fd2 note,anchor=west] at (0.72,-1.38) {$\bar Q$};
      \node[fd2 note,anchor=south,inner sep=1pt] at (0.62,0.44) {$J$（$t=0$）};
      \node[fd2 lab,anchor=north] at (-0.15,-2.15) {$|Q\bar Q,\vec b_\perp\rangle_v$};
      \node[fd2 lab,anchor=west] at (0.8,2.6) {${}_{v'}\langle Q\bar Q,\vec b_\perp|$};
      \node[draw=gray!55,fill=gray!4,rounded corners=2pt,font=\footnotesize,align=left,anchor=north west,inner sep=3pt] at (1.75,-0.55)
        {欧氏表示：\\$\ii\tilde v\cdot D_E=D_\tau-\ii\beta D_z$\\（$\tau$：欧氏时间）};
      \node[fd2 lab] at (0.45,-4.2) {(b)\ HQET 形状因子};
    \fi
  \end{scope}}
  \node[fd2 lab,align=center] at (8.05,-3.1) {$S={}_{v'}\langle Q\bar Q,\vec b_\perp|\,J\,|Q\bar Q,\vec b_\perp\rangle_v$};
\end{tikzpicture}
\caption{离光锥 TMD 软函数作为快速色源对的形状因子（$v=\gamma(1,\beta,\mathbf 0_\perp)$、$v'=\gamma'(1,-\beta',\mathbf 0_\perp)$，$\beta=\tanh Y$、$\beta'=\tanh Y'$，坐标 $(t,z,\vec b_\perp)$，$\vec b_\perp$ 方向按透视画出）。
(a)~两个订书钉形 Wilson 回路（$t<0$ 沿 $v$，$t>0$ 沿 $v'$，在 $-t$ 与 $t'$ 处由横向连接闭合）的时序真空期望，除以 $\sqrt{ZZ}$（$Z$ 为矩形回路）后取 $t,t'\to\infty$；箭头为路径方向：$0$ 处沿世界线指向未来（色源 $Q$），$\vec b_\perp$ 处指向过去（$\bar Q$）。
(b)~等价的 HQET 图像：间隔 $b_\perp$ 的色单态静态 $Q\bar Q$ 对以速度 $v$ 运动，在 $t=0$ 被等时的流 $J$ 踢到 $v'$，软函数是 ${}_{v'}\langle Q\bar Q,\vec b_\perp|J|Q\bar Q,\vec b_\perp\rangle_v$，因而有欧氏表示（\cref{thm:yzliu-soft-hqet}，全阶论证 [A]；非微扰的形状因子形式另需\cref{asm:yzliu-heavy-pair}）。
大快度下 $\ln S=(Y+Y')\tilde K+\mathcal D+\order{\ee^{-(Y+Y')}}$，本征软函数 $S_I=\ee^{-\mathcal D}$（VF-YZL-2）。}
\label{fig:yzliu-soft-hqet}
\end{figure}
%%% END fig:yzliu-soft-hqet

%%% BEGIN fig:yzliu-theta-strip
\begin{figure}[htbp]
\centering
\begin{tikzpicture}[fd,x=1cm,y=0.75cm,
  fd2 ax/.style={lcaxis,gray!60!black},
  fd2 tick/.style={draw=gray!60!black,line width=0.45pt},
  fd2 pt/.style={circle,fill=black,draw=none,inner sep=0pt,minimum size=5.4pt},
  fd2 lab/.style={font=\footnotesize,align=left,anchor=west,inner sep=1.5pt},
  fd2 val/.style={font=\footnotesize,anchor=east,inner sep=1.5pt},
  fd2 note/.style={font=\footnotesize}]
  \def\Yo{3.0}
  % ---- the strip |Im theta| <= pi and its boundary lines ----
  \fill[blue!7] (0,-3.1416) rectangle (11.4,3.1416);
  \draw[blue!50!black,line width=0.75pt] (0,3.1416) -- (11.4,3.1416);
  \draw[blue!50!black,line width=0.75pt] (0,-3.1416) -- (11.4,-3.1416);
  \node[fd2 note,text=blue!45!black,anchor=south east] at (11.4,3.1416) {$\Imag\theta=\pi$};
  \node[fd2 note,text=blue!45!black,anchor=north east] at (11.4,-3.1416) {$\Imag\theta=-\pi$};
  \node[fd2 note,text=blue!45!black,align=center] at (9.35,1.62) {$S(\theta)$ 在 $|\Imag\theta|<\pi$ 中全纯，\\连续到边界};
  % ---- axes ----
  \draw[fd2 ax] (-0.7,0) -- (11.85,0) node[right] {$\Real\theta$};
  \draw[fd2 ax] (0,-4.05) -- (0,4.15) node[left] {$\Imag\theta$};
  \foreach \y/\t in {3.1416/{\pi},1.5708/{\frac{\pi}{2}},-1.5708/{-\frac{\pi}{2}},-3.1416/{-\pi}} {
    \draw[fd2 tick] (-0.07,\y) -- (0.07,\y);
    \node[font=\scriptsize,anchor=east,inner sep=1pt] at (-0.1,\y) {$\t$};}
  \node[font=\scriptsize,anchor=north east,inner sep=1.5pt] at (0,0) {$0$};
  % ---- the line Re theta = Y_0 ----
  \draw[gray!60,line width=0.5pt,dash pattern=on 3pt off 2pt] (\Yo,-3.85) -- (\Yo,3.85);
  \draw[fd2 tick] (\Yo,-0.07) -- (\Yo,0.07);
  \node[font=\scriptsize,anchor=north,inner sep=1.5pt] at (\Yo,-3.85) {$\Real\theta=Y_0$};
  % ---- the six configurations (VF-YZL-4) ----
  \foreach \y in {3.1416,1.5708,0,-1.5708,-3.1416} \node[fd2 pt] at (\Yo,\y) {};
  \node[fd2 val] at (\Yo-0.12,3.1416+0.32) {$Y_0+\ii\pi$};
  \node[fd2 val] at (\Yo-0.12,1.5708+0.32) {$Y_0+\frac{\ii\pi}{2}$};
  \node[fd2 val,fill=blue!7] at (\Yo-0.12,0.36) {$Y_0$};
  \node[fd2 val] at (\Yo-0.12,-1.5708+0.32) {$Y_0-\frac{\ii\pi}{2}$};
  \node[fd2 val] at (\Yo-0.12,-3.1416+0.32) {$Y_0-\ii\pi$};
  \node[fd2 lab] at (\Yo+0.15,3.1416+0.36) {类空/类空，过去/未来};
  \node[fd2 lab] at (\Yo+0.15,1.5708+0.36) {类时/类空，过去/未来};
  \node[fd2 lab] at (\Yo+0.15,0.62) {类空/类空，过去/过去（Collins 的 DY 软因子 $S_{\mathrm{DY}}$）};
  \node[fd2 lab] at (\Yo+0.15,-0.42) {类时/类时，过去/未来（形状因子 $S_{\mathrm{FF}}$）};
  \node[fd2 lab] at (\Yo+0.15,-1.5708+0.36) {类时/类空，过去/过去};
  \node[fd2 lab] at (\Yo+0.15,-3.1416+0.36) {类时/类时，过去/过去};
  % ---- light-cone limit ----
  \draw[red!60!black,line width=0.8pt,-{Stealth[length=5.5pt,width=4.6pt]}] (7.3,-2.05) -- (10.9,-2.05);
  \node[fd2 note,text=red!55!black,align=center,anchor=south] at (9.1,-2.0) {$Y\to\infty$：光锥极限};
  \node[fd2 note,text=red!55!black,align=center,anchor=north] at (9.1,-2.12) {（假设对 $\Imag\theta$ 一致）};
\end{tikzpicture}
\caption{离光锥软函数对双曲角的解析性。单一时序的减除软函数只通过双曲角 $\theta(v_1,v_2)$（$\cosh\theta=v_1\cdot v_2/\sqrt{v_1^2v_2^2}$，$\pm\ii0$ 处方见正文）依赖于两个方向，是带形 $-\pi\le\Imag\theta\le\pi$（阴影，$\Real\theta\ge0$）上同一函数 $S(\theta)$ 的取值；
在固定的 $\Real\theta=Y_0$ 上标出两个方向为类空/类时、过去/未来指向的六种组合所对应的 $\theta$（VF-YZL-4），其中 $\theta=Y_0$ 一点同时代表 Collins 的类空 DY 软因子与类时形状因子软函数，
由此在微扰论各阶 $S_{\mathrm{DY}}=S_{\mathrm{FF}}$（\cref{thm:yzliu-analyticity}，全阶论证 [A]；前提为维数正规化微扰论中的 Wightman 结构，\cref{asm:yzliu-wightman-pt}）。
光锥极限 $Y\to\infty$ 能否越过整个带形延拓，取决于该极限对 $\Imag\theta$ 的一致性；它只在 $\Imag\theta\in\{0,\pm\frac{\pi}{2},\pm\pi\}$ 处得到微扰验证（\cref{asm:yzliu-lightcone-uniform}，部分 [P]）。}
\label{fig:yzliu-theta-strip}
\end{figure}
%%% END fig:yzliu-theta-strip

%%% BEGIN fig:yzliu-threshold
\begin{figure}[htbp]
\centering
\begin{tikzpicture}[fd,
  fd2 ax/.style={lcaxis,gray!60!black},
  fd2 soft/.style={gluon,draw=green!45!black},
  fd2 hadron/.style={bound state,line cap=butt},
  fd2 hv/.style={hard blob,circle,minimum size=0.42cm,inner sep=0pt},
  fd2 lab/.style={font=\small},
  fd2 note/.style={font=\footnotesize}]
  % ======================= (a) DIS at x -> 1 =======================
  \begin{scope}
    % remnant and soft radiation between jet and remnant
    \draw[draw=gray!60!black,line width=0.8pt] (1.45,-0.28) -- (4.85,-0.98);
    \draw[fd2 soft] (3.98,1.12) -- (3.82,-0.76);
    \draw[fd2 soft] (4.55,1.46) -- (4.5,-0.9);
    % struck quark, photon, outgoing jet
    \draw[quark=0.5] (1.5,0.1) -- (2.95,0.38);
    \draw[photon] (2.45,1.75) -- (2.95,0.38);
    \draw[quark=0.6] (2.95,0.38) -- (3.6,1.06);
    \node[jetB blob,ellipse,minimum width=1.75cm,minimum height=0.5cm,inner sep=0pt,rotate=40] at (4.12,1.55) {};
    \node[fd2 hv] at (2.95,0.38) {};
    % target
    \draw[fd2 hadron] (-0.75,0) -- (0.3,0);
    \node[hadron blob,ellipse,minimum width=1.3cm,minimum height=0.85cm,inner sep=0pt] at (0.9,0) {$f$};
    % labels
    \node[fd2 note,anchor=south] at (-0.25,0.05) {$P$};
    \node[fd2 note,anchor=north] at (2.15,0.12) {$xP$};
    \node[fd2 note,anchor=east] at (2.45,1.6) {$q$};
    \node[fd2 note,anchor=north west] at (3.02,0.3) {$H$};
    \node[fd2 note,anchor=south west] at (4.6,2.02) {窄喷注 $J$};
    \node[fd2 note,anchor=north,text=green!30!black] at (4.15,-1.08) {软辐射 $S$};
    \node[fd2 note,anchor=north east,text=gray!35!black] at (2.6,-0.55) {残余 $(1-x)P$};
    \node[fd2 note,align=center] at (2.05,-1.95) {$H$：$Q^2$；\ $J$：$p_J^2\sim(1-x)Q^2$；\ $S$：能量 $\sim(1-x)Q$};
    \node[fd2 lab] at (2.05,-2.6) {(a)\ $x\to1$ 的 DIS};
  \end{scope}
  % ======================= (b) quasi-PDF kernel at xi -> 1 =======================
  \begin{scope}[xshift=9.25cm]
    \draw[fd2 ax] (-2.3,1.0) -- (-1.7,1.0) node[right=-1pt,fd2 note] {$z$};
    \draw[fd2 ax] (-2.3,1.0) -- (-2.3,1.6) node[above=-1pt,fd2 note] {$t$};
    % space-like jet function: link + semi-hard gluons (dashed box)
    \draw[region] (-0.95,-1.05) rectangle (3.95,1.05);
    \draw[fd2 soft] (-0.62,-0.62) .. controls (-0.05,-0.82) and (0.85,-0.62) .. (1.15,0);
    \draw[fd2 soft] (1.85,0) .. controls (2.15,0.62) and (3.0,0.82) .. (3.62,0.62);
    \draw[wilson=0.42] (0,0) -- (3.0,0);
    % collinear quark: in along n at 0, out along n at z
    \draw[quark=0.55] (-1.35,-1.35) -- (0,0);
    \draw[quark=0.5] (3.0,0) -- (4.35,1.35);
    \node[fd2 hv] at (0,0) {};
    \node[fd2 hv] at (3.0,0) {};
    % labels
    \node[fd2 note,anchor=south east] at (-1.02,-1.0) {$p$};
    \node[fd2 note,anchor=east] at (-0.25,0.02) {$H$};
    \node[fd2 note,anchor=north west] at (3.1,-0.1) {$H$};
    \node[fd2 note,anchor=south] at (0.0,0.22) {$0$};
    \node[fd2 note,anchor=south] at (3.0,0.2) {$z$};
    \node[fd2 note,anchor=north] at (2.35,-0.14) {$W[z\hat e_z,0]$};
    \node[fd2 note,anchor=south west,text=gray!35!black] at (-0.95,1.07) {虚线框：类空喷注函数 $\tilde J$};
    \node[fd2 note,align=center] at (1.5,-1.95) {$H$：硬，$\zeta_z=4p_z^2$；\ 绿色：半硬胶子，$|1-\xi|\,p^z$};
    \node[fd2 lab] at (1.5,-2.6) {(b)\ $\xi\to1$ 的准 PDF 匹配核};
  \end{scope}
  % ======================= threshold factorization =======================
  \node[fd2 lab] at (6.85,-3.55) {$C\Big(\xi,\dfrac{p^z}{\mu}\Big)\Big|_{\xi\to1}=H\Big(\dfrac{4p_z^2}{\mu^2}\Big)\,p^z\,J_f\Big(\dfrac{(1-\xi)p^z}{\mu},\dfrac{4p_z^2}{\mu^2}\Big)\big[1+\order{1-\xi}\big]$，\quad $J_f=\tilde J\,\ee^{\ii\,\mathrm{sgn}(z)A}$};
\end{tikzpicture}
\caption{阈值因子化的时空图像（示意）。
(a)~$x\to1$ 的 DIS：被击出的夸克在硬顶点 $H$ 吸收虚光子后形成窄喷注 $J$，喷注与靶残余之间只有软辐射 $S$。
(b)~$\xi\to1$ 的准 PDF 匹配核（$\Gamma=\gamma^0$，$\MSbar$，非单态）：共线夸克（动量 $p$，沿类光方向；细线箭头为费米子数流向）在类空 Wilson 线 $W[z\hat e_z,0]$（双线，箭头为路径方向）的两个连接点处发生硬相互作用（灰色顶点 $H$，标度 $\zeta_z=4p_z^2$），并与静态规范线交换半硬胶子（绿色，标度 $|1-\xi|p^z$）；虚线框内的部分组成类空喷注函数 $\tilde J$，乘上相位因子后为 $J_f$。
$H$ 与准 TMD 因子化中的硬核相同（\cref{thm:yzliu-threshold}，全阶论证 [A]；VF-LAM-14，VF-YZL-9）。
这一因子化形式在单圈精确成立；精确的两圈核在阈值极限下与它的全部 41 个奇异项相容（其中两个常数由比较定出）。}
\label{fig:yzliu-threshold}
\end{figure}
%%% END fig:yzliu-threshold
